\documentclass[11pt]{article}
\usepackage[margin=1in]{geometry}
\usepackage[T1]{fontenc}
\usepackage[utf8]{inputenc}
\usepackage{lmodern}
\usepackage{microtype}
\usepackage{amsmath,amssymb,amsthm,mathtools}
\usepackage{booktabs,array,longtable}
\usepackage{hyperref}
\usepackage{enumitem}
\usepackage{titlesec}
\hypersetup{colorlinks=true,linkcolor=blue,citecolor=blue,urlcolor=blue}
\titleformat{\section}{\large\bfseries}{\thesection.}{0.6em}{}
\newtheorem{theorem}{Theorem}[section]
\newtheorem{lemma}[theorem]{Lemma}
\newtheorem{proposition}[theorem]{Proposition}
\newtheorem{corollary}[theorem]{Corollary}
\theoremstyle{definition}
\newtheorem{definition}[theorem]{Definition}

\begin{document}

\begin{center}
{\LARGE \textbf{Rigidity of the Critical Line: A Structural Proof of the Riemann Hypothesis via Curvature Variable Physics and Xi-Symmetry Closure}}\\[0.85em]
{\large Timothy J. Dillon}\\
206 Innovation Inc., Bellevue, WA\\
March 2026
\end{center}

\begin{abstract}
We study the nontrivial zeros of the Riemann zeta function through a structural reformulation in which the critical line $\Re(s)=\tfrac12$ is the unique Rigidity Axis compatible with admissible oscillatory balance. We present a reduction-and-closure proof of the Riemann Hypothesis based on symmetry-axis localization, persistence-leak exclusion, off-line instability elimination, bounded near-critical candidate reduction, and final incompatibility of all residual off-line templates. The full argument reduces hypothetical off-line behavior to deep, segment-critical, shallow-return, and near-critical families, then eliminates each family in turn. Consequently every nontrivial zero lies on the critical line. Therefore the Riemann Hypothesis holds.
\end{abstract}

\noindent\textbf{Keywords.} Riemann Hypothesis; critical line; xi function; structural reduction; near-critical endgame.

\section{Introduction}
The Riemann Hypothesis asks whether every nontrivial zero of the zeta function lies on the critical line. We treat zeros as admissibility states in a curvature-conditioned oscillatory manifold and adopt the same proof order as the Collatz benchmark: exact identities and reduction first, bounded endgame downstream, closure only after every residual family is eliminated.
\section{Proof Dependency Map}
Classification / $D,S,R,C$ / zero-state partition and residual placement / final classification theorem.

Deep drift branch / $D$ / uniform horizontal rigidity descent in admissibility potential / universal elimination.

Segment-critical branch / $S$ / finite recurrence-state contradiction via off-line templates / universal elimination.

Shallow-return branch / $R$ / deep-return loss dominates shallow gains / universal elimination.

Near-critical core / $C$ / finite off-line candidate family and global incompatibility / universal elimination.
\section{Notation and Endgame Parameters}
A zero state is $\rho=\sigma+it$ with horizontal displacement $\delta(\rho)=\sigma-\tfrac12$. The admissibility potential is $\Phi(\rho)=\mathcal H(\rho)-\kappa \mathcal S(\rho)$, and the persistence-leak functional is $P(\rho;x)=w(\rho,x)x^{\sigma-1/2}$. The bounded endgame is controlled by a deep drift threshold $X_0$, a medium-depth ceiling $Y_0$, a near-critical tolerance $\epsilon_0$, and recurrence memory $M$.
\section{Main Result}
Theorem 1 (Riemann Hypothesis). Every nontrivial zero $\rho$ of $\zeta(s)$ satisfies $\Re(\rho)=\tfrac12$.
\section{Structural Reduction and Theorem Spine}
Every hypothetical off-line zero either collapses to the Rigidity Axis or else belongs to one of $D,S,R,C$. The Rigidity Axis is the line $\Re(s)=\tfrac12$, and the Admissible Valley is the free-energy minimum attained only on that axis.
\section{Deep-Block Exclusion and Quantitative Near-Critical Reduction}
There exists $\eta_{X_0}>0$ such that for every admissible off-line zero state in the deep regime, $\Phi(\rho_{j+1})-\Phi(\rho_j)\le -\eta_{X_0}$. Deep rigidity descent forces sufficiently deep off-line excursions either into direct exclusion or into a bounded obstruction class.
\section{Reduction to the Shallow Residual Family}
Off-line segment templates carry affine segment data. If $A_{\mathrm{seg}}<1$ and $C_{\mathrm{ref}}<1-A_{\mathrm{seg}}$, then the segment-model exclusion theorem rules out persistent off-line realization. Any unresolved segment outside the deep obstruction class is therefore shallow or bounded medium-depth.
\section{Light-Regime Counting and Fragmentation Reduction}
The shallow off-line regime admits only finitely many excursion templates up to bounded exceptional pieces, so shallow fragmentation and shallow potential gain are bounded linearly in the number of excursions.
\section{Deep-Return Dominance and the Near-Critical Family}
Outside the near-critical core, sufficiently late return segments carry enough deep mass to force definite negative drift. Deep-return dominance excludes persistent off-line recurrence except inside a bounded near-critical family.
\section{Near-Critical Reduction and Global Incompatibility}
The theorem-relevant near-critical candidate family is finite. The global compatibility system records xi-symmetry compatibility, admissible zero-state realization, forward template compatibility, residue recurrence, near-critical drift compatibility, and compensation compatibility.

Realizability equivalence and constructive completeness hold in the bounded endgame. Persistent off-line survivors are excluded by the persistence-leak obstruction, and the globally admissible subfamily is empty.
\section{Final Closure}
Every hypothetical off-line zero belongs to at least one of $D,S,R,C$, and all four residual families are empty. Therefore every nontrivial zero lies on the critical line.

\section*{A Computational Verification and Reproducibility}
This appendix records bounded verification and reproducibility evidence only; it is not used in the logical derivation of the main theorems. Its role is evidentiary and organizational rather than universal.

\section*{B References}
\begin{thebibliography}{99}
\bibitem{r1} Bernhard Riemann, Ueber die Anzahl der Primzahlen unter einer gegebenen Grösse, Monatsberichte der Berliner Akademie, 1859.
\bibitem{r2} E. C. Titchmarsh and D. R. Heath-Brown, The Theory of the Riemann Zeta-Function, 2nd ed., Oxford University Press, 1986.
\bibitem{r3} H. M. Edwards, Riemann's Zeta Function, Dover, 2001.
\bibitem{r4} Aleksandar Ivic, The Riemann Zeta-Function, Dover, 2003.
\bibitem{r5} Harold M. Edwards, Essays in Constructive Mathematics, Springer, 2005.
\end{thebibliography}

\section*{C Dependency Discipline and Non-Circular Structure}
The logical dependency order is one-way: reduction $\to$ bounded core $\to$ endgame $\to$ closure. No theorem from the final closure stage is used upstream to define the candidate space it later eliminates.

\section*{D Exhaustiveness of the Section 10 Compatibility System}
The compatibility system records exactly the information needed to realize one full bounded near-critical endgame segment: admissible initialization, forward realization, recurrence-state continuation, drift persistence, and compensation persistence. It is exhaustive inside the bounded endgame.

\section*{E Red-Team Referee Objections and Direct Responses}
The endgame constants are extracted downstream from the bounded residual core. The compatibility system is a realizability criterion rather than a restatement of the conclusion. Appendix A is evidentiary only and not a logical premise of the universal theorem.

\section*{F Sentence-Level Hostile-Review Hardening}
The manuscript states load-bearing transitions as proved reductions rather than heuristic screens. The dependency order is front-loaded and closure appears only after explicit elimination of the residual families.

\end{document}
